In accounting and finance, disclosure has long been viewed not as a neutral release of information, but as an equilibrium object shaped by information asymmetry, proprietary costs, reputational concerns, and the strategic incentives surrounding communication (Healy \& Palepu, 2001; Verrecchia, 2001; Beyer et al., 2010). That logic becomes especially important when the underlying activity is difficult to verify in real time. Signaling models make the point formally: communication is most informative when higher-quality senders can distinguish themselves through statements or actions that lower-quality senders cannot cheaply mimic (Spence, 1973; Seidmann \& Winter, 1997). In settings where implementation is costly but language is inexpensive, disclosure can remain informative without cleanly mapping into realized action. It is precisely in such settings that the distinction between substantive and rhetorical disclosure becomes economically meaningful.

Empirical work on washing sharpens that intuition. Studies of greenwashing show that public claims about sustainability can diverge from realized behavior when firms face strong incentives to appear aligned with socially valued goals (Delmas \& Burbano, 2011; Kim \& Lyon, 2015; Lyon \& Montgomery, 2015). Marquis et al.~(2016) document how selective disclosure can produce the appearance of transparency even when the underlying record is weaker, and Hummel and Schlick (2016) show that disclosure quality rather than disclosure quantity governs whether reporting tracks actual performance. Baker et al.~(2024) bring the same logic to diversity-related disclosure and show that public claims can outpace subsequent workforce outcomes. These studies clearly define the economic problem . When a topic is salient, externally valued, and only imperfectly verifiable, firms may have incentives to use language more aggressively than the underlying implementation record would justify.

A parallel line of work now studies AI communication more directly. Basnet et al.~(2025) show that investors respond differently to actionable and speculative AI language in 10-K business descriptions, while Cao et al.~(2024) show that AI disclosure contains forward-looking information about firm outcomes. Bandyopadhyay et al.~(2023) connect a broader AI narrative measure to subsequent stock mispricing, particularly outside the largest firms, and recent working-paper evidence studies divergence between AI talk and AI walk more explicitly (Li, 2026; Donelson et al., 2025; Barrios et al., 2025). The literature is therefore moving beyond raw mention counts and toward the credibility of AI-related language itself. However, the empirical design choices differ. Some studies focus on manual classification, some on broader narrative exposure, and some on talk--walk divergence measured outside mandatory annual filings. What remains comparatively underdeveloped is a filing-based framework that can be audited, scaled, and reused while preserving the distinction between more credible and less credible AI disclosure inside annual 10-Ks.

Measurement is material since the textual-analysis literature has repeatedly shown that corporate language is informative only when measured in a way that fits the institutional setting. Forward-looking language in filings predicts later fundamentals (Li, 2010), earnings-press-release language contains information beyond the quantitative surprise itself (Davis et al., 2012), and textual measures in finance are highly sensitive to how the underlying vocabulary is constructed (Loughran \& McDonald, 2011, 2016). At the same time, tone and wording can be managed strategically rather than passively reflecting fundamentals (Huang et al., 2014). The implication is not that narrative disclosure is unreliable. The empirical content of narrative disclosure depends on whether the measure isolates the relevant dimension of the language rather than collapsing everything into a broad mention count. That point is especially important in an AI setting, where generic references, aspirational language, and implemented uses can all appear in the same filing while conveying very different information.

The innovation literature points in the same direction. Bellstam et al.~(2021) show that richer textual measures capture meaningful differences in firm-level innovation and forecast subsequent patenting and performance, while Kim and Valentine (2023) show that public disclosure improves how markets interpret innovation-related assets. Those papers justify the two sides of the current design. Filing language can contain information about later technological activity. However, not all ways of describing that language will be equally informative. The purpose of this paper hence is to identify whether the composition of that language isolates a subset of disclosures that are systematically less aligned with contemporaneous implementation and subsequent technological realization. That is the narrower question around which this paper is organized.